\section{Margin-Crossing Mechanism}
\label{sec:margin-mechanism}

This section answers: for deterministic $k$-NN, when does a perturbation
change the prediction?

Once the perturbation type is fixed, a deterministic $k$-NN prediction can only
change for one structural reason: the signed class vote inside the ordered
top-$k$ neighborhood crosses zero.

\subsection{Ordered top-$k$ neighborhoods}

For an ordered sample $S$ and query $x$, let
$\pi_S^x(1),\ldots,\pi_S^x(n)$ be the sample indices sorted by increasing
distance to $x$ and then by sample index.  Write
\[
N_k(S,x) = \big(\pi_S^x(1),\ldots,\pi_S^x(k)\big)
\]
for the ordered top-$k$ neighbor list.

\subsection{Signed vote margin}

Define the signed vote margin
\[
\M_k(S,x) = \sum_{r=1}^{k} \big(2y_{\pi_S^x(r)}-1\big).
\]

Each term $(2y-1)$ maps label $0$ to $-1$ and label $1$ to $+1$.  Under the
label-order tie-break $0 \prec 1$, the deterministic classifier predicts label
$1$ exactly when $\M_k(S,x) > 0$, and label $0$ otherwise.

\subsection{Margin-crossing proposition}

\begin{proposition}[Margin crossing criterion]
\label{prop:5.1}
Fix $S$, $k$, and a query $x$.  Let $S'$ be any perturbed sample of the same
ambient metric space, interpreted with the same deterministic tie-breaking
rule.  Then
\[
h_S^{(k)}(x) \neq h_{S'}^{(k)}(x)
\]
if and only if the signed vote margins satisfy
\[
\big(\M_k(S,x) \le 0 < \M_k(S',x)\big)
\quad\text{or}\quad
\big(\M_k(S',x) \le 0 < \M_k(S,x)\big).
\]
\end{proposition}

\begin{proof}
By definition of the deterministic vote rule, the prediction is $1$ exactly for
strictly positive signed margin and $0$ otherwise.  Therefore the prediction
changes precisely when one margin is positive and the other is non-positive.
\end{proof}

\Cref{prop:5.1} is elementary but useful: every perturbation argument can be
reduced to tracking how the ordered top-$k$ membership changes and how the
corresponding signed vote changes.

Two immediate consequences are worth isolating.  First, deleting an occurrence
outside $N_k(S,x)$ cannot change the prediction at $x$ unless the deletion also
shifts the boundary and pulls a new occurrence into the top-$k$ set.  Second, a
replace-one perturbation can act in two different ways: it can change the
geometry of who enters the neighborhood, or it can leave the geometry fixed and
still flip the vote by changing the label attached to a strategically placed
occurrence.

\subsection{Interpretation for the separation}

The margin criterion reveals why LOO and replace-one can give opposite answers.
LOO perturbations can be arranged so that the margin at every deleted
occurrence remains non-positive (same side of the threshold), while a
replacement at a designated query changes one vote from $-1$ to $+1$,
increasing the margin by $2$ and crossing the threshold.  This is the unifying
mechanism behind all constructions in \Cref{sec:minimal-1nn,sec:odd-k-separation,sec:even-k-separation}.

These mechanisms are especially sensitive in the presence of duplicate
occurrences and equal-distance contenders, because the top-$k$ neighborhood is
then controlled by the deterministic tie-breaking convention rather than by
distance alone.  This is by design: tie-breaking is part of the object being
studied, not a nuisance to be eliminated.
